
% !TeX root = _main.tex
\begin{table*}[t]
    \centering
    \caption{登場キャラクター一覧}
    \label{tab:characters}
    \begin{adjustbox}{max width=\textwidth}
        \begin{tabular}{lllcll}
            \hline
            ID & キャラクター      & 種別    & 初出イベントID & 役柄                         & 初出時の状況                             \\
            \hline
            C1 & ハチドリ（主人公）   & 鳥     & 8        & 物語の中心．墜落→回復→他者との関係を通じた心理変化 & 主人公中心視点に切り替わった直後の鳴き声．ドライ音質で自己存在を明示 \\
            C2 & ハチドリの群れ（仲間） & 鳥（複数） & 2--6     & 世界観・混沌状況の提示，主人公の孤立の対照      & 嵐の中で騒ぐ群れの鳴き声（俯瞰視点・リバーブ付き）          \\
            C3 & ウサギ         & 哺乳類   & 18--19   & 主人公の救助者・同行者・安定した判断主体       & 足音のみで接近（正体不明の存在として提示）              \\
            C4 & シカ          & 哺乳類   & 60       & 導き手（川への案内），状況転換の触媒         & 接近する大型動物の足音（主体未確定）                 \\
            C5 & サル          & 哺乳類   & 117      & 誤解→理解という印象転換を担う存在          & 草の揺れによる存在予兆（正体未確定・警戒対象）            \\
            \hline
        \end{tabular}
    \end{adjustbox}
\end{table*}